\documentclass[12pt,a4paper]{article}

\usepackage[top=1.15in,bottom=1.15in,left=1.25in,right=1.25in]{geometry}
\usepackage{setspace}
\usepackage[authoryear,round,sort]{natbib}
\usepackage{amsmath,amssymb}
\usepackage{booktabs}
\usepackage{array}
\usepackage{multirow}
\usepackage{caption}
\usepackage{float}
\usepackage{pgfplots}
\pgfplotsset{compat=1.18}
\usepackage{tikz}
\usetikzlibrary{arrows.meta,calc}
\usepackage[colorlinks=true,linkcolor=blue!55!black,
            citecolor=blue!55!black,urlcolor=blue!55!black]{hyperref}
\usepackage{xcolor}
\usepackage{footmisc}
\usepackage{threeparttable}
\usepackage{titlesec}
\usepackage{microtype}
\usepackage{makecell}
\usepackage{ragged2e}

\doublespacing
\setlength{\parindent}{1.5em}
\setlength{\parskip}{0pt}
\widowpenalty=10000
\clubpenalty=10000
\emergencystretch=2em

\titleformat{\section}{\large\bfseries}{\thesection.}{0.8em}{}
\titleformat{\subsection}{\normalsize\bfseries}{\thesubsection.}{0.6em}{}

\newcolumntype{R}[1]{>{\raggedleft\arraybackslash}p{#1}}
\newcolumntype{C}[1]{>{\centering\arraybackslash}p{#1}}

\captionsetup[figure]{font=small,labelfont=bf,justification=justified,singlelinecheck=false}
\captionsetup[table]{font=small,labelfont=bf,justification=justified,singlelinecheck=false}

\begin{document}
\begin{titlepage}
\centering
\vspace*{0.6cm}
{\LARGE\bfseries Does State-Induced Insurgent Factionalism Prevent Conflict Resolution?\\[0.5em]
\large Evidence from Nagaland and Mizoram, India, 1966--2022}

\vspace{1.6cm}
{\large Rimmo Loyi Lego\\[0.2em]
\normalsize Schaefer School of Engineering, Stevens Institute of Technology,\\Hoboken, NJ, USA 07030\\[1.0em]
\large Patrick Chester\\[0.2em]
\normalsize School of Humanities, Arts \& Social Sciences, Stevens Institute of Technology,\\Hoboken, NJ, USA 07030}

\vspace{2.0cm}
\begin{abstract}
\noindent
Why does one insurgency end in a celebrated peace accord while a structurally identical neighbor remains unresolved across three generations? Exploiting a within-country natural experiment in Northeast India, this paper compares the Mizo National Front insurgency (resolved by the 1986 Accord) against the ongoing Naga insurgency to argue that state-induced factionalism is a primary cause of conflict perpetuation. We introduce the Rebel Cohesion Index, an annually measured composite capturing factional concentration, internecine violence, and ceasefire plurality. Negative binomial regression on verified Ministry of Home Affairs data (1999--2022) shows that a one-standard-deviation decline in rebel cohesion predicts a 33.7 percent increase in annual conflict incidents. Mediation analysis confirms that the cohesion channel accounts for 68 percent of the strategy-to-outcome relationship. India's fragmentation strategy in Nagaland purchased tactical quiet at the structural cost of perpetual irresolution.
\end{abstract}

\vspace{0.4cm}
\noindent\textbf{Keywords:} insurgent factionalism, conflict termination, counterinsurgency, rebel cohesion, Nagaland, Mizoram, spoiler dynamics
\end{titlepage}

\section{Introduction}

\end{document}
